\emph{Setting.} For $2N$ generators $I^{2} = (-1)^{N(2N-1)}(-1)^{q} = (-1)^{N+q}$, so $I^{2} = -1$ exactly when $N + q$ is odd; $I$ anticommutes with every $e_{i}$ and commutes with $E$. By clause (ii) of \claimref{synthesis.signed-directions-and-the-integer-clifford-algebra}, $e_{S}e_{T} = \pm e_{S\triangle T}$, $e_{S}^{2} = \pm 1$ and $e_{S}e_{T} = (-1)^{b(S,T)}e_{T}e_{S}$ with $b(S,T) = |S||T| + |S\cap T| \bmod 2$. $Ie_{S} = \pm e_{S^{c}}$ and $(Ie_{S})^{2} = I^{2}e_{S}^{2} = -e_{S}^{2}$, so of $e_{S}$ and $Ie_{S}$ exactly one has square $+1$. For even $S$, $T$ the form $b(S,T) = |S\cap T| \bmod 2$ is alternating, the full set is in its radical, and by clause (iii) of that theorem the radical of $b$ on all subsets is trivial for an even number of generators, so the radical on even subsets (the orthogonal complement of the full set) is the full set alone: $b$ is nondegenerate on $V$. The $2D^{2} = 2^{2N-1}$ elements $h_{c}$, $Ih_{c}$ are the even blades up to sign, a basis of $E$. If $b(c,d) = 0$, $(h_{c}h_{d})^{2} = 1$, so $h_{c}h_{d} = \pm h_{c+d}$; if $b(c,d) = 1$, $(h_{c}h_{d})^{2} = -1$, so $h_{c}h_{d} = \pm Ih_{c+d}$. Hence $\langle h_{c}h_{d}\rangle_{0} = \delta_{cd}$, and $x_{c} = D\langle\rho h_{c}\rangle_{0}$.

\emph{(i).} Put $k = Ih_{c}$, $k^{2} = -1$. If $h_{d}$ commutes with $h_{c}$ it commutes with $R_{c}$. If it anticommutes, $R_{c}h_{d}R_{c}^{-1} = \tfrac12(1 + k)h_{d}(1 - k) = \tfrac12 h_{d}(1 - k)^{2} = -h_{d}k = -Ih_{d}h_{c}$, which is $\pm h_{c+d}$ by the setting. So conjugation by $R_{c}$ permutes the set $\{\pm h_{d}\}$ and fixes $I$; an automorphism in $\Gamma$ is determined by its action on that finite set, so $\Gamma$ is finite. $R_{c}^{2} = 2k$ and $I$ is central in $E$, so conjugation by $h_{c}$ lies in $\Gamma$; it multiplies $h_{d}$ by $(-1)^{b(c,d)}$. An automorphism carries idempotents to idempotents and $\Gamma$ preserves the real span of the $h_{c}$ and the scalar part, so states go to states. Transitivity. Write $t_{c}$ for the transvection of $c$. If $c \ne c'$ are nonzero and $b(c,c') = 1$, $t_{c+c'}(c) = c'$. If $b(c,c') = 0$, the functionals $b(\cdot,c)$, $b(\cdot,c')$ are independent, so there is $w$ with $b(w,c) = b(w,c') = 1$, and $t_{w+c'}t_{c+w}(c) = c'$. For pairs: given $(c,d)$ and $(c',d')$ with $b = 0$, first move $c$ to $c'$; then $d$, $d'$ lie in $c'^{\perp}\setminus\{0, c'\}$ and it suffices to move $d$ to $d'$ by transvections $t_{v}$ with $b(v,c') = 0$, which fix $c'$. If $b(d,d') = 1$, $v = d + d'$ does it. If $b(d,d') = 0$ and $d \ne d'$, choose $w$ with $b(w,c') = 0$, $b(w,d) = b(w,d') = 1$: the three functionals are independent unless $d' = d + c'$, in which case the third condition follows from the first two and $c'$, $d$ are independent; then $t_{w+d'}t_{d+w}$ carries $d$ to $d'$ and both factors fix $c'$. The orders $24$ and $11520$ are computed by closing the signed permutations (check). A generator with $k_{c}$ a bivector is a rotor; for $k_{c}$ of grade four the reverse fixes $k_{c}$ and $R_{c}\tilde R_{c} = (1 + k_{c})^{2} = 2k_{c}$; the orbits of the bivector generators alone in $\Cl(7,1)$ and $\Cl(10,0)$ are computed in the check.

\emph{(ii).} The scalar part of $\rho^{2} = \rho$ is $D^{-2}(1 + \sum_{c\ne0}x_{c}^{2}) = D^{-1}$. The coefficient of $h_{u}$, $u \ne 0$, in $\rho^{2}$ is $D^{-2}\bigl(2x_{u} + \sum_{c+d=u}x_{c}x_{d}\,[h_{c}h_{d}]_{h_{u}}\bigr)$ over nonzero $c$, $d$; pairs with $b(c,d) = 1$ contribute to $Ih_{u}$ only, and there $(c,d)$ and $(d,c)$ cancel, as they must since $\rho$ has no $Ih_{u}$ part. Equating with $x_{u}/D$ gives the second identity. The pairs: $c \in u^{\perp}\setminus\{0, u\}$, and $|u^{\perp}| = \tfrac12D^{2}$.

\emph{(iii).} Write $\langle\cdot\rangle$ for $\langle\cdot\rangle_{\Gamma}$; it is invariant under replacing $\rho$ by $g\rho g^{-1}$, $g \in \Gamma$. The coordinates of $g\rho g^{-1}$ are $x_{c}(g\rho g^{-1}) = D\langle\rho\,g^{-1}(h_{c})\rangle_{0}$. \emph{Sign flips.} For a monomial $x_{c_{1}}\cdots x_{c_{j}}$ with $c_{1} + \dots + c_{j} \ne 0$ there is $u$ with $b(u, c_{1} + \dots + c_{j}) = 1$, and conjugation by $h_{u}$ multiplies the monomial by $-1$: its mean is $0$. This gives $\langle x_{c}\rangle = 0$, $\langle x_{c}x_{d}\rangle = 0$ for $c \ne d$, and $\langle x_{c}x_{d}x_{e}\rangle = 0$ unless $e = c + d$ (in particular it vanishes when two of the three classes coincide, the third being nonzero). \emph{Second moment.} By transitivity on nonzero classes $\langle x_{c}^{2}\rangle = a$ does not depend on $c$ (a $g^{-1}$ carrying $h_{c}$ to $\pm h_{c'}$ gives $x_{c}^{2}(g\rho g^{-1}) = x_{c'}^{2}(\rho)$), and the first identity of (ii) gives $(D^{2} - 1)a = D - 1$, $a = 1/(D+1)$. \emph{Third moment, anticommuting.} Let $e = c + d$, $b(c,d) = 1$; then $b(c,e) = b(d,e) = 1$. Conjugation by $R_{e}$ sends $h_{c} \mapsto -Ih_{c}h_{e} = sh_{d}$, $h_{d} \mapsto -Ih_{d}h_{e} = s'h_{c}$ and fixes $h_{e}$; multiplying, $ss'h_{d}h_{c} = I^{2}h_{c}h_{e}h_{d}h_{e} = h_{c}h_{d} = -h_{d}h_{c}$, so $ss' = -1$ and $x_{c}x_{d}x_{e}$ changes sign: its mean is $0$. \emph{Third moment, commuting.} Let $b(c,d) = 0$, $c \ne d$, $e = c + d$ and $T(c,d) = \epsilon(c,d)x_{c}x_{d}x_{c+d}$. If $g^{-1}(h_{c}) = \sigma h_{c'}$ and $g^{-1}(h_{d}) = \tau h_{d'}$ then $g^{-1}(h_{c+d}) = \epsilon(c,d)\sigma\tau\,h_{c'}h_{d'} = \epsilon(c,d)\epsilon(c',d')\sigma\tau\,h_{c'+d'}$, so $T(c,d)(g\rho g^{-1}) = T(c',d')(\rho)$. By transitivity on commuting pairs $\langle T(c,d)\rangle = \kappa$ for all of them. Multiply the second identity of (ii) by $x_{u}$ and average: $(\tfrac12D^{2} - 2)\kappa = (D - 2)a$, so for $D > 2$, $\kappa = 2/((D+1)(D+2))$; for $D = 2$ there is no commuting pair. All values are independent of $\rho$.

\emph{Corollary to (iii).} For a proper subset $A$ of $2a$ generators no even subset of $A$ has its complement in $A$, so the $2^{2a-1} = \tfrac12d_{A}^{2}$ even subsets are distinct classes, the empty one being $0$; the restriction of a state to the algebra they span has purity $d_{A}^{-1}\sum x_{c}^{2}$ in the trace of the $d_{A}$-dimensional module of $A$, and (iii) gives the mean $d_{A}^{-1}(1 + (\tfrac12d_{A}^{2} - 1)/(D+1)) = (D + \tfrac12d_{A}^{2})/(d_{A}(D+1)) = (d_{B} + d_{A})/(d_{A}d_{B} + 2)$ with $D = \tfrac12d_{A}d_{B}$.

\emph{(iv).} With $I^{2} = -1$ and $I$ central, $E$ is a central simple algebra over $\mathrm{span}\{1, I\}$ of dimension $D^{2}$, the $D\times D$ matrices; the adjoint is the real-linear map fixing every $h_{c}$ and negating every $Ih_{c}$; it is an anti-automorphism because $h_{c}h_{d} = \pm h_{c+d}$ for commuting and $\pm Ih_{c+d}$ for anticommuting units, and $\langle xx^{\dagger}\rangle_{0}$ is the sum of the squared coefficients of $x$; the trace is $D$ times the part along $1$ and $I$. \emph{Invariant measures.} Let $\mu$ be a probability measure on states invariant under $\Gamma$ and $f$ a monomial of degree at most three. Then $\int f\,d\mu = \int\langle f\rangle_{\Gamma}\,d\mu$ by invariance, and $\langle f\rangle_{\Gamma}$ is by (iii) a constant, so $\int f\,d\mu$ is that constant. The unitarily invariant measure is invariant under conjugation by every invertible element that preserves the adjoint up to scale, the $R_{c}$ among them ($R_{c}R_{c}^{\dagger} = (1 + k_{c})(1 - k_{c}) = 2$). What follows compares with the imported formula as a cross-check. The formula of Fact \claimref{synthesis.imported-haar-moments-of-random-pure-state} is the sum over permutations of the products of cycle traces, divided by $D(D+1)\cdots(D+n-1)$. For two factors $\bigl(\operatorname{Tr}h_{c}\operatorname{Tr}h_{d} + \operatorname{Tr}h_{c}h_{d}\bigr)/(D(D+1)) = \delta_{cd}/(D+1)$. For three factors every term of the sum over $S_{3}$ containing the trace of a single unit vanishes, the transpositions give $\operatorname{Tr}(h_{c}h_{d})\operatorname{Tr}h_{e} = 0$, and the two three-cycles give $\operatorname{Tr}(h_{c}h_{d}h_{e}) + \operatorname{Tr}(h_{c}h_{e}h_{d})$, which is $2\epsilon(c,d)D$ for a commuting pair with $e = c + d$, zero for an anticommuting pair with $e = c + d$ (the two orders differ by a sign) and zero otherwise; divided by $D(D+1)(D+2)$ these are the values of (iii). A polynomial of degree $n$ in the entries of $\ket{\psi}\bra{\psi}$ is a combination of products of $n$ numbers $\operatorname{Tr}(\ket{\psi}\bra{\psi}h_{c})$, $c \in V$, because the $h_{c}$ span the matrices; so equality on monomials of degree at most three, including those with $x_{0} = 1$, is the design property of order three. The check recomputes both sides independently.

\emph{(v).} A subspace $W$ on which $b$ is nondegenerate has a symplectic basis; the corresponding units generate $r$ mutually commuting copies of the algebra of two anticommuting units of square one, the $2\times2$ matrices over $\mathrm{span}\{1, I\}$, so the $h_{c}$, $c \in W$, span the $m\times m$ matrices, a factor with commutant the $n\times n$ matrices. The reduced state of $\rho$ on it is $m^{-1}\sum_{c\in W}x_{c}h_{c}$ (it has the same expectation on every $h_{c}$, $c \in W$), with purity $m^{-1}\sum_{c\in W}x_{c}^{2}$. By (iii), $\langle P_{W}\rangle = m^{-1}(1 + (m^{2} - 1)/(D + 1)) = (D + m^{2})/(m(D + 1)) = (m + n)/(mn + 1)$ with $D = mn$. This is the value that Fact \claimref{synthesis.imported-haar-moments-of-random-pure-state} states at $n = 2$, at $(d_{A}, d_{B}) = (m, n)$. The even blades of a proper subset $A$ of the generators, $|A| = 2a$, are $2^{2a-1}$ distinct classes, half of $m^{2}$ for $m = 2^{a}$, and do not span a matrix factor.
